\documentclass[a4paper,12pt,oneside]{article}

\usepackage[T2A]{fontenc}
\usepackage[utf8]{inputenc}
\usepackage[russian]{babel}

\usepackage{lmodern}
\usepackage{microtype}

\usepackage{amsmath,amssymb,amsthm,mathtools}

\usepackage{geometry}
\geometry{top=25mm,left=25mm,right=20mm,bottom=25mm}

\usepackage{graphicx}

\usepackage{enumitem}
\setlist[enumerate]{leftmargin=*, itemsep=0.35\baselineskip, topsep=0.35\baselineskip}
\setlist[itemize]{leftmargin=*, itemsep=0.2\baselineskip, topsep=0.25\baselineskip}

\newlist{alphlist}{enumerate}{1}
\setlist[alphlist]{label=\textbf{\alph*)}, leftmargin=*, align=left}

\newlist{numlist}{enumerate}{1}
\setlist[numlist]{label=\textbf{\arabic*)}, leftmargin=*, align=left}

\newcommand{\Task}[1]{%
  \subsection*{#1}%
  \vspace{-0.2\baselineskip}%
}

\setlength{\parindent}{1.25em}
\setlength{\parskip}{0pt}

\begin{document}

\begin{center}
    {\Large\bfseries Карты Карно и минимизация}\\[2mm]
\end{center}
\vspace{2mm}
\hrule
\vspace{4mm}

\Task{Задание 1}
Нахождение СДНФ и СКНФ при помощи карт Карно.
\begin{numlist}
    \item $f=A\overline{B}+B\overline{D}+\overline{A}C$
    \item $f=BD+AC+CD$
    \item $f=\overline{D}+\overline{A}C+\overline{B}D$
    \item $f=\overline{C}+BD+\overline{A}\,\overline{B}$
    \item $f=ABC+BD$
    \item $f=A\overline{C}+AC+BD$
\end{numlist}

\Task{Задание 2}
Нанесение функции на карту Карно:
\[
f=\overline{C}+\overline{A}B+A\overline{B},\qquad
f=\overline{A}B+BCD+\overline{A}\,\overline{C}
\]

\Task{Задание 3}
Алгебраическое упрощение:
\[
f=A\overline{C}+BC+\overline{A}B
\]

\Task{Задание 4}
Импликанты. Чему равна длина вектора минтермов функции, если число её импликант равно:
\[
1)\ 512 \qquad 2)\ 8 \qquad 3)\ 1
\]

\Task{Задание 5}
Методом Квайна найдите сокращённую ДНФ. Укажите число простых импликант, из которых состоит сокращённая форма, и число вхождений переменных:
\[
f(A,B,C)=(0,3,5,6,7)
\]
Что делать, если необходимо найти сокращённую КНФ?

\Task{Задание 6}
Функция задана списком номеров минтермов. Найти минимальную ДНФ и КНФ функции:
\[
f(A,B,C,D)=(0,1,2,6,7,8,12,15)
\]
Может ли у функции не быть ядерных импликант?
\[
\includegraphics[width=0.8\linewidth]{21.1.png}
\]
\[
\includegraphics[width=0.3\linewidth]{21.2.png}
\qquad
\includegraphics[width=0.3\linewidth]{21.3.png}
\]

\end{document}
